\section{Policy Implications and Program Design}
\label{sec:implications}

\subsection{Funding Mechanisms}

The managed freight lane program requires dedicated freight revenue beyond what the Highway Trust Fund can support at current fuel tax rates. Three complementary funding mechanisms are viable at the program scale required:

\textbf{Transponder-based tolling.} Enrolled freight carriers pay a per-mile toll for managed lane access. At \$0.05 per mile --- below the carrier's freight delay savings of \$0.08--0.12 per mile from delay reduction alone --- universal carrier enrollment is economically rational. A \$0.05/mile toll on the managed lanes of all 7 corridors, at 45,600 commercial vehicle equivalents per day and 70\% utilization, generates approximately \$2.3 billion per year in dedicated lane revenue across the portfolio. This covers approximately 18\% of the annualized capital cost of the full program --- meaningful but not sufficient as the sole funding source. Enrollment requires FMCSA database integration (carriers already registered under USDOT number) and transponder deployment at access ramps.

\textbf{CMAQ federal funding.} The Congestion Mitigation and Air Quality Improvement program under IIJA \citep{IIJA2021} provides federal funding for investments that reduce vehicle emissions in non-attainment areas. Managed freight lanes on T1 corridors in urban non-attainment areas (I-95 Northeast corridor, I-80 Bay Area, I-10 Houston, I-75 Atlanta) qualify for CMAQ funding by two mechanisms: (1) diverting freight from general purpose lanes reduces stop-and-go emissions in the GP lanes; (2) platooning reduces fuel consumption and tailpipe emissions on the managed lanes. CMAQ eligibility provides a non-toll federal revenue stream that does not require congressional appropriation beyond the IIJA authorization.

\textbf{Public-private partnership (P3) investment.} The self-funding mechanism of platooning savings creates an investment case for freight carrier capital participation. A carrier operating 500 trucks per day on I-95 managed lanes realizes \$14.6 million per year in fuel savings from platooning at 25\% drag reduction and 30\% fleet penetration. A 10-year concession payment of \$10 million annually --- contributing \$100 million toward I-95 capital over the concession term --- is economically rational for a carrier of that scale. Aggregate P3 participation from major freight carriers (Top 10 LTL carriers operating on T1 corridors) could contribute \$3--5 billion toward Phase 1 capital, reducing the public funding requirement proportionally.

\subsection{Carrier Enrollment and FMCSA Integration}

Managed lane access control requires positive identification of enrolled freight-classified vehicles. The FMCSA Licensing and Insurance database already maintains carrier registration under USDOT number, with vehicle identification through the FMCSA Motor Carrier Management Information System (MCMIS). Managed lane access control can use existing USDOT number transponder credentials: carriers operating under active FMCSA authority and current insurance certification receive automatic transponder enrollment eligibility. Access ramps install RFID transponder readers and license plate recognition cameras for dual-redundant vehicle identification.

This architecture avoids the need for a separate freight vehicle registration system and leverages the existing FMCSA regulatory infrastructure. The access control system enforces not just vehicle class (freight-classified) but also carrier compliance status: carriers with active FMCSA safety fitness ratings of ``Unfit'' or ``Conditional'' (approximately 2.3\% of registered carriers) are excluded from managed lane access during their compliance remediation period, creating a safety benefit alongside the capacity benefit.

\subsection{Investment Phasing Recommendation}

Phase 1 (2025--2030): I-95, I-85, I-35. These three corridors have the highest B/C ratios (2.4:1 to 3.1:1) and the shortest breakeven periods (11--13 years). Combined capital cost: \$31 billion. Phase 1 should be advanced as a package under a single federal program office to enable design standardization, procurement consolidation, and shared enforcement infrastructure. The I-95/I-85 eastern corridor connection is particularly important: a freight carrier operating from New England to the Carolinas traverses both corridors, and coordinated Phase 1 delivery maximizes the shipper SLA benefit for the southeastern manufacturing and distribution sector.

Phase 2 (2030--2035): I-75, I-10. The southeastern and southern transcontinental corridors. Combined capital cost: \$40 billion. I-75 and I-10 share high V/C at their critical segments (Atlanta and Houston, respectively) and strong absolute NPV (\$16B and \$19B). Phase 2 investment should coordinate with the diamond interchange investment at I-75/I-85 Atlanta and I-10/I-95 Jacksonville (see Paper B.4 \citep{ROUTE_B4}): managed lane termini should be designed to connect with diamond zone access ramps, ensuring freight vehicles can transition between the I-75 managed lanes and the I-85 managed lanes through the Atlanta diamond zone without leaving the access-controlled environment.

Phase 3 (2035--2045): I-80, I-90. The northern transcontinental corridors with the longest breakeven periods and lowest B/C ratios. Combined capital cost: \$50 billion. The urban segments of I-80 (Bay Area, Chicago) and I-90 (Chicago, Boston) are the binding constraints on both corridors; phased development beginning with urban segments and deferring rural construction captures the majority of the NPV at lower early capital outlay. At 10\% discount rate, rural segments of both corridors should be evaluated independently on 5-year review cycles as freight volume growth improves their investment case.

\subsection{Interaction with Resilience Investments}

The managed lane program addresses throughput on linear corridor segments; it does not address resilience to closure events. Two interaction effects with existing ROUTE resilience programs are material:

\textbf{Interaction with compound exposure corridors (B.3).} Managed freight lanes improve throughput on I-80 and I-90 but do not reduce vulnerability to Donner Pass or Snoqualmie Pass closures \citep{ROUTE_B3}. A winter storm that closes the Donner Pass segment of I-80 closes the managed lanes along with the general purpose lanes. The managed lane investment is not a substitute for resilience hardening (the Donner freight tunnel proposal); the two programs address different failure modes. The design recommendation is that managed lane construction on closure-vulnerable segments include emergency ramp infrastructure that enables faster deployment of oversize-load alternate routing protocols during closure events.

\textbf{Interaction with diamond interchange zones (B.4).} Managed lane programs should terminate at diamond interchange zone entry points, with dedicated freight-to-freight ramps connecting managed lane termini on one T1 corridor to managed lane starting points on an adjacent T1 \citep{ROUTE_B4}. Without this design coordination, a freight vehicle exiting the I-75 managed lanes at Atlanta must traverse the general purpose I-285 beltway to access I-85 --- losing the PTI benefit of managed lane access precisely at the highest-volume interchange in the Southeast. The design contract for Phase 1 managed lanes on I-75 and I-85 should include the diamond zone connector ramp design as an integrated deliverable, not a separate project.

\subsection{The I-40 Exception as Methodology Validation}

The exclusion of I-40 from the managed lane program is a methodological result, not a policy choice. The ROUTE tier classification assigns I-40 T1 status based on its strategic importance as the southern transcontinental corridor connecting California ports, New Mexico border crossings, and the Southeast \citep{ROUTE_A1}. Tier status does not imply managed lane investment; it implies that the corridor is subject to the full ROUTE analysis pipeline. The capacity analysis in Section~\ref{sec:throughput} determines, independently of tier classification, that I-40 does not require managed lanes at current or projected 10-year demand levels.

The convergence of HPMS V/C analysis (peak V/C = 0.84, below the I2.0 threshold) and ATRI revealed-preference data (zero top-50 bottleneck locations on I-40 \citep{ROUTE_B2}) is a methodological validation of the ROUTE analytical framework: two independent data sources from different methodological traditions reach the same conclusion about the same corridor. This convergence confirms that the managed lane investment program is driven by freight performance data, not by the tier classification assigned to corridors. The tier determines which corridors receive analysis; the analysis determines which corridors receive investment.
